% ============================================================
%  Chapitre 5 — Sprint 3 : Entretiens, Formations & Recommandations
% ============================================================

\chapter*{}
\thispagestyle{empty}

\begin{center}
  \vspace*{2cm}
  {\Huge\bfseries\color{blue!70} Chapitre 5}\\[0.8cm]
  {\LARGE\bfseries\color{blue!70} Sprint 3 : Entretiens, Formations \&}\\[0.3cm]
  {\LARGE\bfseries\color{blue!70} Recommandations IA}\\[1.5cm]

  \begin{tcolorbox}[
    enhanced,
    colback=blue!5,
    colframe=blue!70,
    boxrule=1.5pt,
    arc=4pt,
    width=0.85\textwidth,
    halign=center
  ]
    \large
    Ce chapitre décrit le troisième sprint du projet SIRH, dédié à la gestion
    du cycle de vie des entretiens de sélection. Il couvre l'affectation des
    managers par le Responsable RH, la planification et la conduite des entretiens,
    ainsi que la recommandation intelligente de formations pour les candidats
    nécessitant un renforcement de compétences, grâce à l'intégration de
    l'intelligence artificielle via l'API Gemini de Google.
  \end{tcolorbox}
\end{center}

\addcontentsline{toc}{chapter}{Chapitre 5 : Sprint 3 — Entretiens, Formations \& Recommandations}
\newpage

\renewcommand{\thesection}{\Roman{section}.}
\renewcommand{\thesubsection}{\Roman{section}.\arabic{subsection}.}

% ============================================================
\section{Spécification fonctionnelle du Sprint 3}
% ============================================================

\subsection{Backlog du Sprint 3}

Le Sprint 3 regroupe quatre User Stories portant sur la gestion des entretiens
de sélection et la recommandation de parcours de formation adaptés aux profils
des candidats.

\begin{longtable}{|>{\centering\arraybackslash}p{1.2cm}|>{\arraybackslash}p{2.8cm}|>{\arraybackslash}p{4.5cm}|>{\centering\arraybackslash}p{1.2cm}|>{\arraybackslash}p{4.8cm}|}
  \hline
  \textbf{Id US} &
  \textbf{User Story} &
  \textbf{Description} &
  \textbf{Id tâche} &
  \textbf{Tâche} \\
  \hline
  \endfirsthead

  \hline
  \textbf{Id US} &
  \textbf{User Story} &
  \textbf{Description} &
  \textbf{Id tâche} &
  \textbf{Tâche} \\
  \hline
  \endhead

  \hline
  \endfoot

  \hline
  \endlastfoot

  % ----- US15 -----
  \multirow{3}{*}{\color{blue!70}\textbf{US15}} &
  \multirow{3}{2.8cm}{\color{blue!70} Assigner un manager à l'entretien} &
  \multirow{3}{4.5cm}{En tant que Responsable RH, je veux affecter un manager
  à une candidature retenue afin d'organiser un entretien de sélection.} &
  \color{blue!70} 15.1 & Élaborer les diagrammes du scénario « Assigner un manager ». \\ \cline{4-5}
  & & & \color{blue!70} 15.2 & Implémenter le cas d'utilisation « Assigner un manager ». \\ \cline{4-5}
  & & & \color{blue!70} 15.3 & Tester le cas d'utilisation « Assigner un manager ». \\ \hline

  % ----- US16 -----
  \multirow{3}{*}{\color{blue!70}\textbf{US16}} &
  \multirow{3}{2.8cm}{\color{blue!70} Planifier et conduire l'entretien} &
  \multirow{3}{4.5cm}{En tant que Manager, je veux consulter mes entretiens
  planifiés, envoyer une convocation et enregistrer le résultat afin de statuer
  sur les candidats.} &
  \color{blue!70} 16.1 & Élaborer les diagrammes du scénario « Conduire l'entretien ». \\ \cline{4-5}
  & & & \color{blue!70} 16.2 & Implémenter le cas d'utilisation « Conduire l'entretien ». \\ \cline{4-5}
  & & & \color{blue!70} 16.3 & Tester le cas d'utilisation « Conduire l'entretien ». \\ \hline

  % ----- US17 -----
  \multirow{3}{*}{\color{blue!70}\textbf{US17}} &
  \multirow{3}{2.8cm}{\color{blue!70} Recommander une formation via IA} &
  \multirow{3}{4.5cm}{En tant que Manager, je veux obtenir des recommandations
  d'organismes de formation adaptées aux compétences manquantes du candidat,
  générées par intelligence artificielle.} &
  \color{blue!70} 17.1 & Élaborer les diagrammes du scénario « Recommander une formation ». \\ \cline{4-5}
  & & & \color{blue!70} 17.2 & Implémenter le cas d'utilisation « Recommander une formation ». \\ \cline{4-5}
  & & & \color{blue!70} 17.3 & Tester le cas d'utilisation « Recommander une formation ». \\ \hline

  % ----- US18 -----
  \multirow{3}{*}{\color{blue!70}\textbf{US18}} &
  \multirow{3}{2.8cm}{\color{blue!70} Notifier les acteurs par email} &
  \multirow{3}{4.5cm}{En tant que Système, je veux envoyer automatiquement
  des convocations à l'entretien et des notifications de proposition de formation
  aux employés concernés.} &
  \color{blue!70} 18.1 & Élaborer les diagrammes du scénario « Envoyer les notifications ». \\ \cline{4-5}
  & & & \color{blue!70} 18.2 & Implémenter le cas d'utilisation « Envoyer les notifications ». \\ \cline{4-5}
  & & & \color{blue!70} 18.3 & Tester le cas d'utilisation « Envoyer les notifications ». \\ \hline

\end{longtable}

\subsection{Acteurs impliqués}

\begin{longtable}{|>{\bfseries\arraybackslash}p{3.5cm}|>{\arraybackslash}p{11.5cm}|}
  \hline
  \textbf{Acteur} &
  \textbf{Rôle dans le Sprint 3} \\
  \hline
  \endfirsthead
  \hline
  \textbf{Acteur} &
  \textbf{Rôle dans le Sprint 3} \\
  \hline
  \endhead
  \hline
  \endfoot

  Administrateur RH &
  Consulte les candidatures retenues, désigne le manager responsable de
  l'entretien et déclenche l'envoi de la convocation au candidat. \\ \hline

  Manager &
  Prend en charge les entretiens assignés, consulte le profil du candidat,
  conduit l'entretien et formule sa décision. En cas d'orientation en formation,
  il initie la recommandation par intelligence artificielle et soumet
  la proposition à l'employé. \\ \hline

  Employé &
  Reçoit la convocation à l'entretien par email. Il est ensuite informé de la
  décision et peut accepter ou décliner la proposition de formation depuis
  son espace personnel. \\ \hline

  API Gemini (Google) &
  Service d'intelligence artificielle sollicité par la plateforme pour analyser
  l'écart de compétences du candidat et proposer les trois organismes de formation
  les mieux adaptés à son profil. \\ \hline

\end{longtable}

% ============================================================
\section{Diagramme de cas d'utilisation du Sprint 3}
% ============================================================

Le diagramme suivant illustre l'ensemble des cas d'utilisation du Sprint 3.
Trois acteurs principaux interagissent avec le système : l'Administrateur RH,
le Manager et l'Employé. L'API Gemini intervient en tant qu'acteur secondaire
lors de la phase de recommandation de formation.

\begin{figure}[H]
  \centering
  \includegraphics[width=0.92\textwidth]{images/sprint3_usecase.png}
  \caption{Diagramme de cas d'utilisation — Sprint 3}
  \label{fig:sprint3_usecase}
\end{figure}

% ============================================================
\section{Analyse des cas d'utilisation}
% ============================================================

\subsection{Cas d'utilisation : Assigner un manager à l'entretien}

\begin{longtable}{|>{\bfseries\arraybackslash}p{3.5cm}|>{\arraybackslash}p{11.5cm}|}
  \hline
  \rowcolor{headerblue}
  \multicolumn{2}{|l|}{\textcolor{white}{\textbf{UC — Assigner un manager à l'entretien}}} \\
  \hline
  \endfirsthead
  \hline
  \endhead
  \hline
  \endfoot

  Acteur principal & Administrateur RH \\ \hline

  Pré-conditions &
  L'Administrateur RH est connecté à la plateforme. Une candidature retenue
  est en attente d'affectation. Au moins un manager est disponible dans le
  système. \\ \hline

  Post-conditions &
  Un entretien est planifié et associé au manager désigné. Une convocation
  est envoyée automatiquement au candidat. Le statut de la candidature est
  mis à jour. \\ \hline

  Scénario principal &
  \begin{enumerate}
    \item L'Administrateur RH consulte la liste des candidatures en attente
          d'entretien.
    \item Il sélectionne une candidature et ouvre le formulaire d'affectation.
    \item Le système affiche la liste des managers disponibles.
    \item L'administrateur choisit un manager et renseigne la date, l'heure
          et le lieu de l'entretien.
    \item Il confirme. Le système enregistre l'entretien, envoie une convocation
          par email au candidat et met à jour le statut de la candidature.
  \end{enumerate} \\ \hline

  Scénario alternatif &
  \textbf{[Aucun manager disponible]} : le système affiche un message
  d'avertissement et l'affectation est suspendue jusqu'à l'ajout d'un manager. \\ \hline
\end{longtable}

\begin{figure}[H]
  \centering
  \includegraphics[width=0.88\textwidth]{images/seq_sys_assigner.png}
  \caption{Diagramme de séquence système — Assigner un manager}
  \label{fig:seq_sys_assigner}
\end{figure}

\begin{figure}[H]
  \centering
  \includegraphics[width=\textwidth]{images/seq_det_assigner.png}
  \caption{Diagramme de séquence détaillé — Assigner un manager}
  \label{fig:seq_det_assigner}
\end{figure}

\subsection{Cas d'utilisation : Planifier et conduire l'entretien}

\begin{longtable}{|>{\bfseries\arraybackslash}p{3.5cm}|>{\arraybackslash}p{11.5cm}|}
  \hline
  \rowcolor{headerblue}
  \multicolumn{2}{|l|}{\textcolor{white}{\textbf{UC — Planifier et conduire l'entretien}}} \\
  \hline
  \endfirsthead
  \hline
  \endhead
  \hline
  \endfoot

  Acteur principal & Manager \\ \hline

  Pré-conditions &
  Le Manager est connecté à la plateforme. Un entretien lui a été assigné
  et le candidat a reçu sa convocation. \\ \hline

  Post-conditions &
  La décision de l'entretien est enregistrée. Le statut de la candidature
  est mis à jour et le candidat est notifié du résultat. \\ \hline

  Scénario principal &
  \begin{enumerate}
    \item Le Manager accède à son tableau de bord et consulte ses entretiens
          planifiés.
    \item Il sélectionne un entretien et prend connaissance du profil et
          du score de compatibilité du candidat.
    \item Il conduit l'entretien et saisit ses observations.
    \item Il enregistre sa décision : Retenu, Refusé ou Orientation en formation.
    \item Le système met à jour le statut de la candidature et notifie
          le candidat par email.
  \end{enumerate} \\ \hline

  Scénario alternatif &
  \textbf{[Orientation en formation]} : le manager choisit d'orienter le
  candidat vers une formation. Le système calcule l'écart de compétences
  et sollicite l'API Gemini pour proposer les organismes les plus adaptés.
  Le manager sélectionne un organisme et soumet la proposition à l'employé. \\ \hline
\end{longtable}

\begin{figure}[H]
  \centering
  \includegraphics[width=0.88\textwidth]{images/seq_sys_entretien.png}
  \caption{Diagramme de séquence système — Conduire l'entretien}
  \label{fig:seq_sys_entretien}
\end{figure}

\begin{figure}[H]
  \centering
  \includegraphics[width=\textwidth]{images/seq_det_entretien.png}
  \caption{Diagramme de séquence détaillé — Conduire l'entretien}
  \label{fig:seq_det_entretien}
\end{figure}

\subsection{Cas d'utilisation : Recommander une formation via IA}

\begin{longtable}{|>{\bfseries\arraybackslash}p{3.5cm}|>{\arraybackslash}p{11.5cm}|}
  \hline
  \rowcolor{headerblue}
  \multicolumn{2}{|l|}{\textcolor{white}{\textbf{UC — Recommander une formation via IA}}} \\
  \hline
  \endfirsthead
  \hline
  \endhead
  \hline
  \endfoot

  Acteur principal & Manager (déclenche) / Système IA (traite) \\ \hline

  Pré-conditions &
  L'entretien s'est conclu par une orientation en formation. Le catalogue
  des organismes de formation est disponible dans la plateforme. \\ \hline

  Post-conditions &
  Une proposition de formation est créée et soumise à l'employé.
  L'employé reçoit une notification par email l'invitant à consulter
  et répondre à la proposition. \\ \hline

  Scénario principal &
  \begin{enumerate}
    \item Le Manager déclenche la recommandation depuis le tableau de bord.
    \item Le système identifie les compétences manquantes du candidat par
          rapport aux exigences du poste.
    \item Il transmet ces informations, accompagnées du catalogue des
          organismes, au service d'intelligence artificielle.
    \item L'API Gemini retourne les trois organismes les plus adaptés,
          avec un score et une justification pour chacun.
    \item Le manager consulte les résultats, choisit un organisme et
          confirme la proposition.
    \item La proposition est enregistrée et l'employé est notifié.
  \end{enumerate} \\ \hline

  Scénario alternatif &
  \textbf{[Service IA indisponible]} : le manager peut sélectionner
  manuellement un organisme depuis le catalogue de la plateforme. \\ \hline
\end{longtable}

\begin{figure}[H]
  \centering
  \includegraphics[width=0.88\textwidth]{images/seq_sys_formation_reco.png}
  \caption{Diagramme de séquence système — Recommandation de formation}
  \label{fig:seq_sys_formation_reco}
\end{figure}

\begin{figure}[H]
  \centering
  \includegraphics[width=\textwidth]{images/seq_det_formation_reco.png}
  \caption{Diagramme de séquence détaillé — Recommandation de formation}
  \label{fig:seq_det_formation_reco}
\end{figure}

% ============================================================
\section{Diagramme de classes du Sprint 3}
% ============================================================

Le diagramme de classes ci-dessous présente les entités introduites durant
le Sprint 3 ainsi que leurs relations avec les éléments existants du système.

Les trois nouvelles entités sont :
\begin{itemize}
  \item \textbf{Entretien} : regroupe les informations relatives à un entretien
        de sélection — le manager responsable, la date, le lieu et la décision
        finale. Elle est associée à une candidature interne ou à un candidat externe.
  \item \textbf{Organisme de formation} : représente un prestataire de formation
        référencé dans le catalogue de la plateforme, avec ses domaines
        de compétences, son mode de délivrance et son statut de partenariat.
  \item \textbf{Proposition de formation} : matérialise la proposition soumise
        à l'employé à l'issue de l'entretien, avec le suivi de son état
        (proposée, acceptée, refusée ou en cours).
\end{itemize}

\begin{figure}[H]
  \centering
  \includegraphics[width=\textwidth]{images/class_sprint3.png}
  \caption{Diagramme de classes — Sprint 3}
  \label{fig:class_sprint3}
\end{figure}

% ============================================================
\section{Réalisation}
% ============================================================

\subsection{Interface Responsable RH — Affectation du manager}

L'interface RH permet de sélectionner un manager et de lui affecter une
candidature pour la planification de l'entretien. La convocation est
générée et envoyée automatiquement dès la validation.

\begin{figure}[H]
  \centering
  \includegraphics[width=0.92\textwidth]{images/ihm_assigner_manager.png}
  \caption{Interface RH — Affectation d'un manager à un entretien}
  \label{fig:ihm_assigner_manager}
\end{figure}

\subsection{Interface Manager — Tableau de bord des entretiens}

Le tableau de bord centralise les entretiens assignés au manager. Il lui
permet de consulter le profil de chaque candidat, de saisir ses observations
et d'enregistrer sa décision à l'issue de l'entretien.

\begin{figure}[H]
  \centering
  \includegraphics[width=0.92\textwidth]{images/ihm_manager_interviews.png}
  \caption{Interface Manager — Tableau de bord des entretiens}
  \label{fig:ihm_manager_interviews}
\end{figure}

\subsection{Interface Manager — Recommandations de formation}

Après une décision d'orientation en formation, la plateforme affiche les
trois organismes recommandés par l'intelligence artificielle, avec leur
score de pertinence et une justification. Le manager sélectionne l'organisme
le plus adapté avant de soumettre la proposition.

\begin{figure}[H]
  \centering
  \includegraphics[width=0.92\textwidth]{images/ihm_formation_reco.png}
  \caption{Interface Manager — Recommandations de formation générées par IA}
  \label{fig:ihm_formation_reco}
\end{figure}

\subsection{Interface Employé — Proposition de formation}

L'employé consulte depuis son espace personnel la proposition de formation
qui lui a été soumise. Il peut visualiser l'organisme recommandé et les
compétences ciblées, puis accepter ou décliner la proposition.

\begin{figure}[H]
  \centering
  \includegraphics[width=0.92\textwidth]{images/ihm_employee_formation.png}
  \caption{Interface Employé — Consultation de la proposition de formation}
  \label{fig:ihm_employee_formation}
\end{figure}

% ============================================================
\section*{Conclusion}
\addcontentsline{toc}{section}{Conclusion}
% ============================================================

Ce troisième sprint a permis de mettre en place le cycle complet de gestion
des entretiens au sein de la plateforme SIRH, depuis l'affectation du manager
jusqu'à la soumission d'une proposition de formation personnalisée à l'employé.

L'intégration de l'intelligence artificielle via l'API Gemini constitue la
valeur ajoutée principale de ce sprint : en analysant l'écart entre le profil
du candidat et les exigences du poste, la plateforme génère des recommandations
d'organismes de formation pertinentes et justifiées, réduisant le temps de
décision du manager et assurant des orientations adaptées aux besoins réels.

Le sprint suivant portera sur la gestion des candidats externes et le
développement d'outils de pilotage avancés pour les managers.
